\documentclass[a4paper,12pt]{article}
\usepackage[ngerman]{babel}
\usepackage[utf8]{inputenc}
\usepackage[T1]{fontenc}
\usepackage{enumitem}
\usepackage{amsmath,amssymb}
\usepackage{geometry}
\geometry{left=25mm,right=25mm,top=25mm,bottom=25mm}
\usepackage{parskip}

\usepackage{tikz}
\usetikzlibrary{circuits.ee.IEC, positioning}

\newif\ifshowsolutions
\showsolutionstrue  % comment to hide solutions

\begin{document}

\begin{center}
  {\LARGE\bfseries Einführung in die Elektrotechnik \\[6pt] Testklausur}\\[12pt]
  {\large Besprechung: 12.01.2026}\\[6pt]
  \rule{\textwidth}{0.4pt}
\end{center}

\section{Verständnis: Multiple Choice}

Kreuze jeweils eine Antwort pro Frage an.

\vspace{8pt}

\subsection*{Grundlagen}

\begin{enumerate}[label=\textbf{\arabic*.}, leftmargin=*, itemsep=8pt]

  \item Die Driftbewegung von Elektronen in einem Leiter entsteht hauptsächlich durch:
    \begin{itemize}[label=$\square$]
      \item[A)] Temperaturgradienten
      \item[B)] Magnetfelder
      \item[C)] das elektrische Feld
      \item[D)] mechanische Einwirkung
    \end{itemize}

\ifshowsolutions
Lösung: C
\fi

  \item Bei einem metallischen Leiter führt eine Temperaturerhöhung üblicherweise zu:
    \begin{itemize}[label=$\square$]
      \item[A)] abnehmendem Widerstand
      \item[B)] plötzlichem Isolationsverhalten
      \item[C)] unverändertem Widerstand
      \item[D)] steigendem Widerstand
    \end{itemize}

\ifshowsolutions
Lösung: D
\fi

\end{enumerate}

\subsection*{Gleichstrom}

\begin{enumerate}[label=\textbf{\arabic*.}, resume, leftmargin=*, itemsep=8pt]

  \item Die Knotenregel (erste Kirchhoffsche Regel) besagt:
    \begin{itemize}[label=$\square$]
      \item[A)] Die Summe der Spannungen in einer Masche ist Null.
      \item[B)] Die Summe der Ströme in einem Knoten ist Null.
      \item[C)] Die Summe der Widerstände in einem Knoten ist konstant.
      \item[D)] Leistung wird in einem Knoten gespeichert.
    \end{itemize}

\ifshowsolutions
Lösung: B
\fi

  \item Eine ideale Spannungsquelle hat charakteristisch:
    \begin{itemize}[label=$\square$]
      \item[A)] festen Spannungswert unabhängig vom Strom
      \item[B)] festen Stromwert unabhängig von der Last
      \item[C)] unendlichen Innenwiderstand
      \item[D)] keinen Einfluss auf angeschlossene Schaltungen
    \end{itemize}

\ifshowsolutions
Lösung: A
\fi

  \item Das Knotenpotentialverfahren arbeitet primär mit:
    \begin{itemize}[label=$\square$]
      \item[A)] Maschenströmen
      \item[B)] Knotenspannungen
      \item[C)] Ersatzquellen
      \item[D)] Frequenzanalysen
    \end{itemize}

\ifshowsolutions
Lösung: B
\fi

  \item Für maximale Leistungsübertragung von einer Quelle an eine ohmsche Last muss gelten:
    \begin{itemize}[label=$\square$]
      \item[A)] Lastwiderstand $\gg$ Quellenwiderstand
      \item[B)] Lastwiderstand $\ll$ Quellenwiderstand
      \item[C)] Lastwiderstand gleich Quelleninnenwiderstand
      \item[D)] Lastwiderstand = 0
    \end{itemize}

\ifshowsolutions
Lösung: C
\fi

  \item Ein Amperemeter wird üblicherweise:
    \begin{itemize}[label=$\square$]
      \item[A)] parallel zur Messgröße geschaltet
      \item[B)] in Reihe zur Messgröße geschaltet
      \item[C)] ohne Innenwiderstand betrieben
      \item[D)] nur bei Wechselstrom eingesetzt
    \end{itemize}

\ifshowsolutions
Lösung: B
\fi

\end{enumerate}

\subsection*{Felder \& Bauelemente}

\begin{enumerate}[label=\textbf{\arabic*.}, resume, leftmargin=*, itemsep=8pt]

  \item Ein Kondensator speichert primär:
    \begin{itemize}[label=$\square$]
      \item[A)] magnetische Energie
      \item[B)] mechanische Energie
      \item[C)] thermische Energie
      \item[D)] elektrische Energie
    \end{itemize}

\ifshowsolutions
Lösung: D
\fi

  \item Bei konstanter Gleichspannung verhält sich ein Kondensator (auf lange Sicht) wie:
    \begin{itemize}[label=$\square$]
      \item[A)] Kurzschluss
      \item[B)] Unterbrechung (kein Stromfluss)
      \item[C)] Widerstand mit fixem Wert
      \item[D)] Spannungsquelle
    \end{itemize}

\ifshowsolutions
Lösung: B
\fi

  \item Magnetfelder entstehen insbesondere durch:
    \begin{itemize}[label=$\square$]
      \item[A)] ruhende Ladungen
      \item[B)] bewegte Ladungen bzw. Strom
      \item[C)] Temperaturänderungen
      \item[D)] Lichtstrahlung
    \end{itemize}

\ifshowsolutions
Lösung: B
\fi

  \item Elektrische Induktion tritt auf, wenn:
    \begin{itemize}[label=$\square$]
      \item[A)] die Spannung konstant bleibt
      \item[B)] die Leiterlänge konstant bleibt
      \item[C)] der magnetische Fluss durch eine Fläche sich ändert
      \item[D)] keine Ladungen vorhanden sind
    \end{itemize}

\ifshowsolutions
Lösung: C
\fi

  \item Eine Spule (Induktivität) wirkt gegen schnelle Stromänderungen, weil sie:
    \begin{itemize}[label=$\square$]
      \item[A)] elektrische Energie speichert
      \item[B)] magnetische Energie speichert
      \item[C)] Wärme erzeugt
      \item[D)] Ladungen neutralisiert
    \end{itemize}

\ifshowsolutions
Lösung: B
\fi

  \item Ein idealer Transformator kann:
    \begin{itemize}[label=$\square$]
      \item[A)] Gleichspannung in höhere Gleichspannung umwandeln
      \item[B)] Wechselspannung in andere Wechselspannung (andere Amplitude) umwandeln
      \item[C)] die Frequenz einer Wechselspannung dauerhaft ändern
      \item[D)] Energie erzeugen
    \end{itemize}

\ifshowsolutions
Lösung: B
\fi

  \item Ein transienter Vorgang in einem elektrischen Netzwerk bedeutet:
    \begin{itemize}[label=$\square$]
      \item[A)] sofortige, sprungartige Änderung ohne zeitliche Entwicklung
      \item[B)] zeitlich veränderlicher Einschwingvorgang bis zum stationären Zustand
      \item[C)] permanenter Zustand ohne Änderung
      \item[D)] rein mathematische Betrachtung ohne physikalische Bedeutung
    \end{itemize}

\ifshowsolutions
Lösung: B
\fi

\end{enumerate}

\subsection*{Wechselstrom}

\begin{enumerate}[label=\textbf{\arabic*.}, resume, leftmargin=*, itemsep=8pt]

  \item Ein Phasor (Zeiger) stellt dar:
    \begin{itemize}[label=$\square$]
      \item[A)] die komplette zeitliche Kurve eines Signals
      \item[B)] nur Gleichstromgrößen
      \item[C)] ausschließlich den Realteil einer Größe
      \item[D)] eine komplexe Amplitude (Betrag und Phase) bei fixer Frequenz
    \end{itemize}

\ifshowsolutions
Lösung: D
\fi

  \item Der Blindwiderstand eines Kondensators ist bei steigender Frequenz:
    \begin{itemize}[label=$\square$]
      \item[A)] größer
      \item[B)] kleiner
      \item[C)] unverändert
      \item[D)] unvorhersehbar
    \end{itemize}

\ifshowsolutions
Lösung: B
\fi

  \item Blindleistung entsteht in einem Wechselstromkreis, wenn:
    \begin{itemize}[label=$\square$]
      \item[A)] Spannung und Strom phasenverschoben sind
      \item[B)] Spannung und Strom exakt gleichphasig sind
      \item[C)] keine Energie übertragen wird
      \item[D)] nur Gleichstrom fließt
    \end{itemize}

\ifshowsolutions
Lösung: A
\fi

  \item Ein Hochpass-Filter lässt in erster Näherung durch:
    \begin{itemize}[label=$\square$]
      \item[A)] niedrige Frequenzen
      \item[B)] keine Signale
      \item[C)] ausschließlich Gleichstrom
      \item[D)] hohe Frequenzen
    \end{itemize}

\ifshowsolutions
Lösung: D
\fi

  \item Resonanz in einem Schwingkreis bedeutet typischerweise:
    \begin{itemize}[label=$\square$]
      \item[A)] minimale Impedanz bzw. maximaler Strom (bei Serie)
      \item[B)] maximale Dämpfung
      \item[C)] kein Einfluss der Last
      \item[D)] sofortige Abklingzeit
    \end{itemize}

\ifshowsolutions
Lösung: A
\fi

  \item Ein symmetrisches Drehstromsystem hat als praktischen Vorteil:
    \begin{itemize}[label=$\square$]
      \item[A)] geringere Frequenz
      \item[B)] dass nur zwei Leiter benötigt werden
      \item[C)] die Erzeugung eines rotierenden Magnetfelds (z.\,B. für Motoren)
      \item[D)] kein Ausgleichsstrom möglich
    \end{itemize}

\ifshowsolutions
Lösung: C
\fi

\end{enumerate}

\newpage

\section{Gleichstromnetzwerke}

\subsection{Ersatzwiderstand}

Gegeben ist folgendes Widerstandsnetzwerk mit $R_1 = 10\,\Omega$, $R_2 = 20\,\Omega$, $R_3 = 30\,\Omega$, $R_4 = 15\,\Omega$:

\begin{center}
\begin{tikzpicture}[circuit ee IEC, scale=1]
  \draw
  (0,0) to[resistor={info={$R_1$}}] (2,0)
        to[resistor={info={$R_3$}}] (4,0)
        -- (4,-2)
  (0,-2) to[resistor={info={$R_2$}}] (2,-2)
        -- (4,-2)
  (0,0) -- (0,-2)
  (4,0) to[resistor={info={$R_4$}}] (6,0)
        -- (6,-2) -- (4,-2);
\end{tikzpicture}
\end{center}

Berechnen Sie den Ersatzwiderstand des Netzwerks.

\ifshowsolutions
Lösung:
\[
R_{12} = \frac{R_1 R_2}{R_1 + R_2} = \frac{10 \cdot 20}{30} = 6.67\,\Omega
\]
\[
R_{123} = 6.67 + 30 = 36.67\,\Omega
\]
\[
R = \frac{36.67 \cdot 15}{36.67 + 15} \approx 10.7\,\Omega
\]

\subsection{Netzwerkanalyse}

Gegeben ist folgendes Netzwerk mit $R_1 = 6\,\Omega$, $R_2 = 12\,\Omega$, $R_3 = 4\,\Omega$ und einer Spannungsquelle von $24\,V$:

\begin{center}
\begin{tikzpicture}[circuit ee IEC]
  \draw
  (0,0) to[voltage source={info={$U$}}] (0,3)
        to[resistor={info={$R_1$}}] (3,3)
        -- (3,1.5)
  (3,1.5) to[resistor={info'={$R_2$}}] (6,1.5)
        -- (6,0)
  (3,1.5) to[resistor={info={$R_3$}}] (6,3)
        -- (6,0) -- (0,0);
\end{tikzpicture}
\end{center}

Berechnen Sie

\begin{enumerate}
  \item den Gesamtstrom der Quelle.
  \item die Teilströme durch $R_2$ und $R_3$.
  \item die Leistung an $R_3$.
\end{enumerate}

\ifshowsolutions
Lösung:
\[
R_{23} = \frac{12 \cdot 4}{16} = 3\,\Omega
\]
\[
R = 6 + 3 = 9\,\Omega
\]
\[
I = \frac{24}{9} = 2.67\,\text{A}
\]
\[
U_{23} = I \cdot 3 = 8\,\text{V}
\]
\[
I_2 = \frac{8}{12} = 0.67\,\text{A}, \quad I_3 = \frac{8}{4} = 2\,\text{A}
\]
\[
P_3 = I_3^2 R_3 = 16\,\text{W}
\]
\fi

\subsection{Knotenpotentialverfahren}

Gegeben ist folgendes Gleichstromnetzwerk mit Spannungsquelle $U = 20\,\mathrm{V}$, $R_1 = 5\,\Omega$, $R_2 = 10\,\Omega$, $C = 100\,\mu\mathrm{F}$:

\begin{center}
\begin{tikzpicture}[circuit ee IEC, scale=1]
  \draw
  % Spannungsquelle
  (0,0) to[voltage source={info={$U$}}] (0,4)

  % R1
  (0,4) to[resistor={info'={$R_1$}}] (4,4)

  % Knoten
  (4,4) node[circle,fill,inner sep=2pt] {}

  % R2
  (4,4) to[resistor={info'={$R_2$}}] (8,4)
        -- (8,0)

  % Kondensator
  (4,4) to[capacitor={info'={$C$}}] (4,0)

  % Masse
  (0,0) -- (8,0);
\end{tikzpicture}
\end{center}

Der untere Leiter ist als Bezugspotential (Masse) gewählt und es wird der stationäre Zustand ($t \rightarrow \infty$) betrachtet.

Bestimmen Sie

\begin{enumerate}
  \item das Potential am eingezeichneten Knoten,
  \item die Spannung am Kondensator,
  \item den Strom durch den Kondensator,
  \item die im Kondensator gespeicherte Energie.
\end{enumerate}

\ifshowsolutions
Lösung:
\[
\frac{V_1 - 20}{R_1} + \frac{V_1 - 0}{R_2} = 0
\]
\[
2(V_1 - 20) + V_1 = 0
\]
\[
V_1 = \frac{40}{3} \approx 13.33\,\text{V}
\]
\[
U_C = V_1 - 0 = 13.33\,\text{V}
\]
\[
i_C = C \frac{du_C}{dt} = 0
\]
\[
W = \frac{1}{2} C U_C^2 \approx 8.9\,\text{mJ}
\]
\fi

\newpage

\section{Wechselstromnetzwerke}

\subsection{Ersatzimpedanz}

Gegeben ist folgendes Wechselstromnetzwerk mit $R_1 = 10\,\Omega$, $R_2 = 20\,\Omega$, $L = 50\,\mathrm{mH}$, $C = 100\,\mu\mathrm{F}$, $\omega = 1000\,\mathrm{rad/s}$:

\begin{center}
\begin{tikzpicture}[circuit ee IEC, scale=1]
  \draw
  % Eingang
  (0,0) node[left]{A}
  to[resistor={info'={$R_1$}}] (3,0)

  % Parallelschaltung
  (3,0) -- (3,2)
  to[inductor={info'={$L$}}] (6,2)
  -- (6,0)

  (3,0) -- (3,-2)
  to[capacitor={info'={$C$}}] (6,-2)
  -- (6,0)

  % R2
  to[resistor={info'={$R_2$}}] (9,0)
  node[right]{B};
\end{tikzpicture}
\end{center}

Bestimmen Sie

\begin{enumerate}
  \item die komplexen Impedanzen der einzelnen Bauelemente,
  \item die komplexe Ersatzimpedanz zwischen den Klemmen A und B,
  \item Betrag und Phase der Ersatzimpedanz.
\end{enumerate}

\ifshowsolutions
Lösung:
\[
\underline{Z}_R = R
\]
\[
\underline{Z}_L = j \omega L = j \cdot 1000 \cdot 0.05 = j50\,\Omega
\]
\[
\underline{Z}_C = \frac{1}{j \omega C} = \frac{1}{j \cdot 1000 \cdot 100 \cdot 10^{-6}} = -j10\,\Omega
\]
\[
\underline{Z}_{LC} = \left( \frac{1}{Z_L} + \frac{1}{Z_C} \right)^{-1} = \frac{1}{j0.08} = -j12.5\,\Omega
\]
\[
\underline{Z} = R_1 + Z_{LC} + R_2 = (30 - j12.5)\,\Omega
\]
\[
|\underline{Z}| = \sqrt{30^2 + 12.5^2} \approx 32.5\,\Omega
\]
\[
\varphi = \arctan\left(\frac{-12.5}{30}\right) \approx -22.6^\circ
\]
\fi

\subsection{Komplexe Netzwerkanalyse}

Gegeben sei folgendes Netzwerk mit sinusförmiger Spannungsquelle $U = 230\,\mathrm{V}$, $f = 50\,\mathrm{Hz}$ und den Bauelementen $R = 20\,\Omega$, $L = 100\,\mathrm{mH}$:

\begin{center}
\begin{tikzpicture}[circuit ee IEC]
  \draw
  (0,0) to[voltage source={info={$U$}}] (0,3)
        to[resistor={info={$R$}}] (3,3)
        to[inductor={info={$L$}}] (6,3)
        -- (6,0) -- (0,0);
\end{tikzpicture}
\end{center}

Bestimmen Sie

\begin{enumerate}
  \item die komplexe Impedanz $Z$,
  \item den Effektivstrom $I$,
  \item die Phasenverschiebung zwischen Strom und Spannung,
  \item den Leistungsfaktor $\cos\varphi$,
  \item die Wirkleistung $P$,
  \item die Blindleistung $Q$.
\end{enumerate}

\ifshowsolutions
Lösung:
\[
\omega = 2\pi f = 314\,\text{rad/s}
\]
\[
X_L = \omega L = 31.4\,\Omega
\]
\[
\underline{Z} = 20 + j31.4 \Rightarrow |Z| = 37.3\,\Omega
\]
\[
I = \frac{230}{37.3} = 6.17\,\text{A}
\]
\[
\varphi = \arctan\left(\frac{31.4}{20}\right) = 57.5^\circ
\]
\[
\cos\varphi = 0.54
\]
\[
P = U I \cos\varphi = 766\,\text{W}
\]
\[
Q = U I \sin\varphi = 1180\,\text{var}
\]
\fi

\newpage

\section*{Formelsammlung}

\subsection*{Kondensator}

\[
i_C(t) = C \frac{du_C(t)}{dt}
\]

\[
W = \frac{1}{2} C U^2
\]

\subsection*{Wechselstrom}

\[
\omega = 2 \pi f
\]

\[
\underline{Z}_L = j \omega L
\]

\[
\underline{Z}_C = \frac{1}{j \omega C}
\]

\[
\underline{Z} = R + jX
\]

\[
P = U I \cos \varphi
\]

\[
Q = U I \sin \varphi
\]

\end{document}